% !TEX root = ../main.tex

\chapter{Eigene Klassifikation adaptiver
Lernverfahren}\label{eigene-klassifikation-adaptiver-lernverfahren}

Aufbauend auf den in Kapitel 2 identifizierten Anforderungen wird im
Folgenden eine eigenständige Klassifikation adaptiver Lernverfahren
entwickelt. Ziel ist nicht die Einordnung nach Anwendungskontexten oder
Plattformtypen, sondern die systematische Beschreibung der Verfahren
anhand ihrer funktionalen Rolle innerhalb adaptiver Systeme sowie
weiterer ergänzender Merkmale. Die Klassifikation soll Verfahren
unterschiedlicher Forschungsfelder vergleichbar machen und gleichzeitig
die Einordnung konkreter Algorithmen ermöglichen.

\section{Wahl der
Klassifikationsstruktur}\label{wahl-der-klassifikationsstruktur}

Die in Kapitel 2 analysierte Literatur macht deutlich, dass adaptive
Lernverfahren durch eine hohe methodische Vielfalt gekennzeichnet sind.
Neben probabilistischen, logistischen und tiefen Lernmodellen existieren
Recommender-Systeme, Lernpfadverfahren, psychometrische Ansätze und
Reinforcement-Learning-Verfahren, die teilweise unterschiedliche
Aufgaben innerhalb adaptiver Systeme übernehmen, teilweise jedoch auch
ähnliche Probleme mit unterschiedlichen Modellierungsansätzen lösen.
Gleichzeitig treten viele Verfahren in mehreren Forschungskontexten auf
und werden in realen Systemen häufig miteinander kombiniert. Die
Entwicklung einer eigenen Klassifikation erfordert daher zunächst die
Frage, welche Strukturform für eine systematische Einordnung dieser
Verfahren geeignet ist.

Eine rein hierarchische Taxonomie bietet den Vorteil einer klaren und
intuitiven Struktur. Verfahren können schrittweise in Ober- und
Unterkategorien eingeordnet werden, wodurch sich eine übersichtliche
Gliederung ergibt. Für adaptive Lernverfahren erweist sich ein solcher
Ansatz jedoch als nur eingeschränkt geeignet. Viele Verfahren lassen
sich nicht eindeutig einer einzelnen Kategorie zuordnen, da hybride
Systeme mehrere Modellierungs- und Entscheidungsverfahren kombinieren.
Zudem treten identische Verfahren teilweise in unterschiedlichen
Anwendungskontexten auf. Eine ausschließlich hierarchische Struktur
würde daher entweder Mehrfachzuordnungen erzwingen oder relevante
Unterschiede zwischen Verfahren verdecken.

Demgegenüber ermöglicht eine rein facettierte Taxonomie die unabhängige
Beschreibung verschiedener Eigenschaften eines Verfahrens. Datenbasis,
Modellierungsansatz, Entscheidungsmechanismus oder Interpretierbarkeit
könnten getrennt voneinander betrachtet werden. Dieser Ansatz bietet
eine hohe analytische Flexibilität und eignet sich insbesondere für den
Vergleich heterogener Verfahren. Für die Zielsetzung dieser Arbeit
besitzt eine ausschließlich facettierte Struktur jedoch den Nachteil,
dass sie keinen intuitiven primären Ordnungsrahmen bereitstellt.
Verfahren würden lediglich als Kombination verschiedener Merkmale
erscheinen, ohne dass ihre funktionale Rolle innerhalb adaptiver Systeme
unmittelbar sichtbar wird.

Für die vorliegende Arbeit erscheint daher eine hierarchisch verankerte
facettierte Mehrdimensionaltaxonomie am geeignetsten. Sie kombiniert die
Vorteile beider Ansätze: Einerseits werden Verfahren zunächst anhand
ihrer funktionalen Rolle innerhalb des adaptiven Prozesses eingeordnet.
Andererseits ermöglichen zusätzliche Beschreibungsdimensionen eine
differenzierte Charakterisierung der Verfahren. Auf diese Weise lassen
sich sowohl konkrete Algorithmen als auch hybride Systeme systematisch
erfassen, ohne die Übersichtlichkeit der Klassifikation zu verlieren.

\subsection{Herleitung der
Hauptdimensionen}\label{herleitung-der-hauptdimensionen}

Die Analyse der in Kapitel 2 betrachteten Verfahrensfamilien zeigt, dass
adaptive Lernsysteme trotz ihrer methodischen Vielfalt zwei grundlegende
Aufgaben erfüllen müssen. Zum einen benötigen sie Mechanismen zur
Erfassung und Modellierung des aktuellen Lernzustands. Zum anderen
müssen sie auf Grundlage dieser Informationen Entscheidungen über
Inhalte, Aufgaben, Interventionen oder Lernpfade treffen. Unabhängig
davon, ob ein System probabilistische Modelle, Recommender-Verfahren,
psychometrische Ansätze oder Reinforcement Learning verwendet, lassen
sich die zugrunde liegenden Verfahren letztlich einer dieser beiden
Aufgaben zuordnen.

Verfahren wie Bayesian Knowledge Tracing, Performance Factors Analysis,
Item-Response-Theory-Modelle oder Deep Knowledge Tracing dienen primär
der Schätzung von Wissen, Fähigkeiten oder zukünftigen Leistungen. Ihr
Ziel besteht darin, einen internen Zustand des Lernenden möglichst
präzise zu modellieren. Recommender-Systeme, Lernpfadverfahren,
Computerized Adaptive Testing oder Reinforcement-Learning-Verfahren
treffen demgegenüber Entscheidungen darüber, welche Inhalte, Aufgaben
oder Interventionen als Nächstes ausgewählt werden sollen. Sie greifen
dabei häufig auf Informationen zurück, die zuvor durch Verfahren der
Zustandsmodellierung bereitgestellt wurden.

Aus dieser Beobachtung ergibt sich die funktionale Rolle eines
Verfahrens als geeignete primäre Klassifikationsdimension. Die
vorgeschlagene Taxonomie unterscheidet daher zunächst zwischen Verfahren
der Zustandsmodellierung und Verfahren der Adaptionsentscheidung.

\section{Die vorgeschlagene
Klassifikation}\label{die-vorgeschlagene-klassifikation}

\subsection{Zustandsmodellierung}\label{zustandsmodellierung}

Die erste Hauptklasse der vorgeschlagenen Klassifikation umfasst
Verfahren zur Zustandsmodellierung des Lernenden. Gemeinsam ist diesen
Verfahren, dass ihre unmittelbare Ausgabe nicht in einer konkreten
adaptiven Entscheidung besteht, sondern in der Schätzung eines internen
Zustands des Lernenden. Ziel ist es, Wissen, Fähigkeiten, Kompetenzen,
Präferenzen oder zukünftige Leistungen auf Grundlage beobachtbarer
Interaktionen möglichst präzise zu modellieren. Die erzeugten
Zustandsinformationen dienen anschließend häufig als Grundlage für
weitere adaptive Entscheidungen, etwa bei der Auswahl von
Lernressourcen, der Generierung von Lernpfaden oder der Steuerung
adaptiver Tests.

Die Analyse der in Kapitel 2 betrachteten Forschungsfelder zeigt, dass
Verfahren der Zustandsmodellierung in nahezu allen Bereichen adaptiver
Lernsysteme eine zentrale Rolle einnehmen. Im Knowledge Tracing steht
die Schätzung des Beherrschungsgrads einzelner Wissenskomponenten im
Vordergrund. Im Computerized Adaptive Testing werden Fähigkeiten über
psychometrische Modelle geschätzt, während Recommender-Systeme und
Lernpfadverfahren häufig Informationen über Wissen, Präferenzen oder
Lernfortschritte als Eingabe für spätere Entscheidungen nutzen. Trotz
unterschiedlicher Anwendungskontexte verfolgen diese Verfahren somit
dieselbe grundlegende Funktion: die Modellierung eines internen Zustands
des Lernenden auf Basis beobachtbarer Daten.

Aus der Literatur lassen sich innerhalb dieser Hauptklasse fünf zentrale
Untergruppen ableiten. Diese unterscheiden sich primär hinsichtlich
ihrer Modellierungslogik, ihrer Annahmen über den Lernprozess sowie der
Art und Weise, wie Zustandsinformationen aus den verfügbaren Daten
gewonnen werden. Die vorgeschlagene Klassifikation unterscheidet daher
zwischen psychometrischen Fähigkeitsmodellen, probabilistisch-temporalen
Wissensmodellen, logistischen und online aktualisierten
Performanzmodellen, merkmalsbasierten Machine-Learning-Modellen sowie
tiefen sequenziellen und repräsentationslernenden Modellen.

\subsubsection{Psychometrische
Fähigkeitsmodelle}\label{psychometrische-fuxe4higkeitsmodelle}

Psychometrische Fähigkeitsmodelle bilden die erste Unterklasse der
Zustandsmodellierung. Ihr gemeinsames Ziel besteht darin, latente
Fähigkeiten oder Kompetenzen eines Lernenden auf Grundlage beobachtbarer
Antwortmuster zu schätzen. Im Gegensatz zu vielen später entwickelten
Lernmodellen steht dabei weniger die Modellierung eines zeitlichen
Lernverlaufs als vielmehr die möglichst präzise Messung eines aktuellen
Fähigkeitsniveaus im Vordergrund \citep{Han2018, Liu2024}.

Zu den wichtigsten Vertretern dieser Modellfamilie zählen das
Rasch-Modell sowie die verschiedenen Modelle der Item Response Theory
(IRT). Diese Verfahren beschreiben die Wahrscheinlichkeit einer
korrekten Antwort als Funktion von Personen- und Itemparametern und
ermöglichen dadurch die Schätzung latenter Fähigkeiten aus Testdaten. In
adaptiven Testsystemen dienen sie häufig als Grundlage für die Auswahl
informativer Aufgaben und die fortlaufende Aktualisierung von
Fähigkeitswerten \citep{Han2018, Liu2024}.

Innerhalb der vorgeschlagenen Klassifikation werden psychometrische
Modelle als eigenständige Unterklasse geführt, da ihr primärer Fokus auf
der diagnostischen Fähigkeitsmessung liegt. Während spätere
Knowledge-Tracing-Modelle Lernprozesse über mehrere Interaktionen hinweg
modellieren, konzentrieren sich psychometrische Verfahren auf die
möglichst zuverlässige Schätzung eines latenten Zustands zu einem
gegebenen Zeitpunkt \citep{Shen2024, Han2018}. Typische Vertreter dieser
Unterklasse sind das Rasch-Modell sowie die 2PL- und 3PL-Varianten der
Item Response Theory.

\subsubsection{Probabilistisch-temporale
Wissensmodelle}\label{probabilistisch-temporale-wissensmodelle}

Probabilistisch-temporale Wissensmodelle bilden die zweite Unterklasse
der Zustandsmodellierung. Ihr gemeinsames Merkmal besteht darin, dass
sie Wissen nicht als statischen Zustand, sondern als dynamischen
Lernprozess modellieren. Ziel dieser Verfahren ist die Schätzung der
Wahrscheinlichkeit, mit der ein Lernender eine bestimmte
Wissenskomponente beherrscht, wobei diese Schätzung nach jeder neuen
Lerninteraktion aktualisiert wird \citep{Shen2024, Abdelrahman2023}.

Charakteristisch für diese Modellfamilie ist die explizite Modellierung
von Wissensentwicklung über die Zeit. Viele Verfahren dieser Klasse
basieren auf probabilistischen Zustandsmodellen, insbesondere
Hidden-Markov-Modellen oder Bayes-Netz-artigen Strukturen. Wissen wird
dabei als latenter, nicht direkt beobachtbarer Zustand aufgefasst, der
anhand beobachtbarer Antworten oder Interaktionen inferiert wird
\citep{Shen2024}.

Der wichtigste Vertreter dieser Unterklasse ist Bayesian Knowledge
Tracing (BKT). Das von Corbett und Anderson entwickelte Verfahren
modelliert Lernen als Folge von Zustandsübergängen und berücksichtigt
dabei insbesondere Anfangswissen, Lernwahrscheinlichkeit sowie die
Möglichkeiten von Guessing, also die Wahrscheinlichkeit, dass eine
Aufgabe trotz fehlender Beherrschung korrekt beantwortet wird, und
Slipping, der umgekehrte Fall einer falschen Antwort trotz vorhandenen
Wissens. Aufbauend auf dem klassischen BKT wurden zahlreiche
Erweiterungen entwickelt, die beispielsweise individuelle
Lernendenparameter, Itemschwierigkeiten oder komplexere
Abhängigkeitsstrukturen berücksichtigen
\citep{Shen2024, Abdelrahman2023}.

Innerhalb der vorgeschlagenen Klassifikation werden
probabilistisch-temporale Wissensmodelle als eigenständige Unterklasse
geführt, da sie im Unterschied zu psychometrischen Fähigkeitsmodellen
nicht primär die Messung eines aktuellen Fähigkeitsniveaus verfolgen,
sondern die zeitliche Entwicklung von Wissen modellieren. Typische
Vertreter dieser Modellfamilie sind Bayesian Knowledge Tracing (BKT)
sowie Dynamic Bayesian Knowledge Tracing (DBKT)
\citep{Shen2024, Abdelrahman2023}.

\subsubsection{Logistische und online aktualisierte
Performanzmodelle}\label{logistische-und-online-aktualisierte-performanzmodelle}

Logistische und online aktualisierte Performanzmodelle bilden die dritte
Unterklasse der Zustandsmodellierung. Im Mittelpunkt dieser Verfahren
steht nicht die explizite Modellierung eines verborgenen
Wissenszustands, sondern die Vorhersage zukünftiger Leistungen auf
Grundlage beobachtbarer Lernhistorien. Ziel ist es, die
Wahrscheinlichkeit einer korrekten Antwort aus bisherigen Erfolgen,
Fehlern und Übungserfahrungen abzuleiten \citep{Pelanek2017, Shen2024}.

Charakteristisch für diese Modellfamilie ist die Verwendung
statistischer oder online aktualisierter Modelle zur
Performanzschätzung. Im Gegensatz zu probabilistischen Wissensmodellen
wird Wissen dabei häufig nicht als latenter Zustand modelliert, sondern
indirekt über beobachtbare Leistungsindikatoren beschrieben. Viele
Verfahren dieser Klasse nutzen logistische Regressionen, um den Einfluss
früherer Interaktionen auf zukünftige Antworten zu schätzen
\citep{Pelanek2017, Shen2024}.

Zu den wichtigsten Vertretern dieser Unterklasse zählen Learning Factors
Analysis (LFA) und Performance Factors Analysis (PFA). Beide Verfahren
modellieren die Wahrscheinlichkeit korrekter Antworten auf Basis der
bisherigen Bearbeitungshistorie eines Lernenden. Während LFA
insbesondere die Anzahl früherer Übungsgelegenheiten berücksichtigt,
unterscheidet PFA zusätzlich zwischen erfolgreichen und fehlerhaften
Bearbeitungen und kann dadurch Lernfortschritte differenzierter abbilden
\citep{Pelanek2017, Shen2024}.

Ebenfalls dieser Unterklasse zuzuordnen sind Elo-basierte Verfahren.
Diese interpretieren die Interaktion zwischen Lernendem und Aufgabe als
eine Art Wettbewerb, bei dem sowohl die Fähigkeit des Lernenden als auch
die Schwierigkeit der Aufgabe nach jeder Bearbeitung fortlaufend
aktualisiert werden. Dadurch eignen sich Elo-Verfahren insbesondere für
adaptive Systeme, die schnelle Online-Aktualisierungen bei
vergleichsweise geringem Modellierungsaufwand benötigen
\citep{Pelanek2017, Shen2024}.

Innerhalb der vorgeschlagenen Klassifikation werden logistische und
online aktualisierte Performanzmodelle als eigenständige Unterklasse
geführt, da sie im Unterschied zu probabilistischen Wissensmodellen
keine explizite Repräsentation eines verborgenen Wissenszustands
anstreben, sondern unmittelbar aus beobachtbaren Leistungsdaten auf
zukünftige Performanz schließen. Typische Vertreter dieser Modellfamilie
sind LFA, PFA, PFAE sowie Elo-basierte Verfahren
\citep{Pelanek2017, Shen2024}.

\subsubsection{Merkmalsbasierte
Machine-Learning-Modelle}\label{merkmalsbasierte-machine-learning-modelle}

Merkmalsbasierte Machine-Learning-Modelle bilden die vierte Unterklasse
der Zustandsmodellierung. Im Unterschied zu den zuvor beschriebenen
Modellfamilien beruhen diese Verfahren nicht auf spezifischen Annahmen
über Wissenszustände, Lernprozesse oder Fähigkeitsmodelle. Stattdessen
nutzen sie allgemeine Methoden des Machine Learning, um aus
beobachtbaren Merkmalen Muster zu erkennen und Eigenschaften von
Lernenden vorherzusagen \citep{Wang2025, Orji2022}.

Die Eingabedaten dieser Verfahren können sehr unterschiedlich ausfallen.
Neben Antwortmustern und Leistungsdaten werden häufig Interaktionsdaten,
Nutzungsverhalten, demographische Merkmale oder weitere Informationen
über Lernende berücksichtigt. Ziel ist dabei nicht ausschließlich die
Schätzung von Wissen oder Fähigkeiten, sondern auch die Vorhersage
weiterer Lernendenmerkmale wie Motivation, Engagement, Lernstilen,
Persönlichkeitsmerkmalen oder affektiven Zuständen \citep{Orji2022}.

Zu den typischen Vertretern dieser Unterklasse zählen
Entscheidungsbäume, Random Forests, k-Nearest-Neighbor-Verfahren (kNN),
Support Vector Machines (SVM), Clustering-Verfahren sowie verschiedene
Fuzzy-Logic-Ansätze. Darüber hinaus werden in der Literatur auch
Bayesian Networks und faktorenanalytische Verfahren zur Modellierung von
Lernendenmerkmalen beschrieben \citep{Wang2025, Orji2022}.

Innerhalb der vorgeschlagenen Klassifikation werden diese Verfahren als
eigenständige Unterklasse geführt, da sie nicht auf speziell
entwickelten Wissens- oder Fähigkeitsmodellen beruhen, sondern
allgemeine datengetriebene Lernverfahren zur Vorhersage von
Lernendenmerkmalen einsetzen. Im Unterschied zu tiefen neuronalen
Verfahren arbeiten merkmalsbasierte Machine-Learning-Modelle
typischerweise mit explizit definierten Eingabemerkmalen, die aus den
verfügbaren Lern- und Interaktionsdaten abgeleitet werden. Die
relevanten Eigenschaften der Lernenden werden dabei nicht automatisch
gelernt, sondern dem Modell in Form von Features bereitgestellt. Sie
stellen damit eine flexiblere, aber häufig auch weniger
domänenspezifische Alternative zu psychometrischen, probabilistischen
oder logistischen Modellen dar.

\subsubsection{Tiefe sequenzielle und repräsentationslernende
Modelle}\label{tiefe-sequenzielle-und-repruxe4sentationslernende-modelle}

Die fünfte Unterklasse der Zustandsmodellierung umfasst tiefe
sequenzielle und repräsentationslernende Modelle. Im Unterschied zu
psychometrischen, probabilistischen oder logistischen Ansätzen beruhen
diese Verfahren nicht auf explizit definierten Wissenszuständen oder
fest vorgegebenen Lernparametern. Stattdessen lernen sie latente
Repräsentationen direkt aus den Interaktionsdaten von Lernenden und
nutzen diese zur Vorhersage zukünftiger Leistungen oder Wissenszustände
\citep{Shen2024}.

Besonders verbreitet sind diese Verfahren im Bereich des Knowledge
Tracing. Zu den bekanntesten Vertretern zählen Deep Knowledge Tracing
(DKT), speicherbasierte Modelle sowie aufmerksamkeitsbasierte und
transformerbasierte Knowledge-Tracing-Verfahren \citep{Shen2024}.
Gemeinsam ist diesen Ansätzen, dass sie Interaktionssequenzen als
zeitliche Datenfolgen behandeln und komplexe Abhängigkeiten zwischen
früheren und späteren Lernaktivitäten modellieren können.

Charakteristisch für diese Modellfamilie ist die automatische Bildung
latenter Repräsentationen. Anstatt Wissen explizit über Parameter oder
Wahrscheinlichkeitsmodelle zu beschreiben, lernen die Modelle interne
Zustandsrepräsentationen, die zur Vorhersage zukünftiger Antworten
genutzt werden. Dadurch können auch komplexe und nichtlineare
Zusammenhänge erfasst werden, die in klassischen Modellen häufig nicht
abgebildet werden \citep{Shen2024}.

Die hohe Vorhersageleistung dieser Verfahren geht jedoch häufig mit
einer geringeren Interpretierbarkeit einher. Während psychometrische
oder probabilistische Modelle ihre Annahmen explizit offenlegen, bleiben
die internen Repräsentationen tiefer neuronaler Modelle oftmals schwer
nachvollziehbar. Die Literatur diskutiert deshalb regelmäßig den
Zielkonflikt zwischen Vorhersagegüte und Erklärbarkeit tiefer
Knowledge-Tracing-Modelle \citep{Shen2024}.

Innerhalb der vorgeschlagenen Klassifikation werden tiefe sequenzielle
und repräsentationslernende Modelle als eigenständige Unterklasse
geführt, da sie Wissen nicht über explizite Fähigkeits- oder
Wissensparameter beschreiben, sondern latente Zustände datengetrieben
aus großen Interaktionssequenzen erlernen. Typische Vertreter dieser
Kategorie sind DKT sowie weitere neuronale Sequenz- und
Transformer-Modelle \citep{Shen2024}.

\subsection{Adaptionsentscheidung}\label{adaptionsentscheidung}

Die zweite Hauptklasse der vorgeschlagenen Klassifikation umfasst
Verfahren der Adaptionsentscheidung. Im Unterschied zu den Verfahren der
Zustandsmodellierung besteht ihre unmittelbare Aufgabe nicht in der
Schätzung eines internen Zustands des Lernenden, sondern in der Auswahl,
Steuerung oder Planung einer konkreten Systemaktion. Grundlage dieser
Entscheidungen sind häufig die durch Learner Models,
Knowledge-Tracing-Verfahren oder andere Zustandsmodelle bereitgestellten
Informationen. Das Ergebnis dieser Verfahren ist somit nicht eine
Diagnose des Lernenden, sondern eine Entscheidung darüber, wie das
adaptive System auf den aktuellen Zustand reagieren soll.

Verfahren der Adaptionsentscheidung finden sich in nahezu allen
Bereichen adaptiver Lernsysteme. Computerized-Adaptive-Testing-Systeme
wählen auf Basis der aktuellen Fähigkeitsschätzung die nächste Aufgabe
aus. Recommender-Systeme empfehlen Lernressourcen oder Aufgaben, die zum
Lernenden passen. Lernpfadverfahren planen ganze Sequenzen von
Lernaktivitäten, während Reinforcement-Learning-Verfahren langfristige
Strategien zur Optimierung des Lernerfolgs entwickeln. Trotz
unterschiedlicher Anwendungskontexte verfolgen diese Ansätze damit
dieselbe grundlegende Funktion: die Auswahl einer geeigneten Folgeaktion
auf Grundlage verfügbarer Informationen über den Lernenden und die
Lernumgebung.

Die in Kapitel 2 analysierte Literatur zeigt zugleich, dass sich
adaptive Entscheidungen auf sehr unterschiedliche Weise realisieren
lassen. Einige Verfahren beruhen auf explizit definierten Regeln oder
Expertenwissen, andere nutzen informationsbasierte Auswahlkriterien,
Empfehlungssysteme oder Optimierungsverfahren. Darüber hinaus
unterscheiden sich die Verfahren hinsichtlich ihres Planungshorizonts.
Während manche Ansätze lediglich die nächste Aufgabe oder Ressource
auswählen, planen andere vollständige Lernpfade oder optimieren
langfristige Entscheidungsstrategien über viele Interaktionen hinweg.

Aus der Literatur lassen sich innerhalb dieser Hauptklasse fünf zentrale
Untergruppen ableiten. Diese unterscheiden sich primär hinsichtlich
ihres Entscheidungsmechanismus sowie der Art der erzeugten
Adaptionsentscheidung. Die vorgeschlagene Klassifikation unterscheidet
daher zwischen regel- und wissensbasierten Adaptionsverfahren,
informationsmaximierenden Selektionsverfahren, Verfahren der
Ressourcenempfehlung und Aufgabenwahl, Verfahren zur Lernpfadgenerierung
und Sequenzplanung sowie Verfahren der sequentiellen Optimierung unter
Unsicherheit.

\subsubsection{Regel- und wissenbasierte
Adaptionsverfahren}\label{regel--und-wissenbasierte-adaptionsverfahren}

Regel- und wissensbasierte Adaptionsverfahren bilden die erste
Unterklasse der Adaptionsentscheidung. Gemeinsam ist diesen Verfahren,
dass adaptive Entscheidungen nicht durch statistische Optimierung oder
datengetriebene Lernverfahren erzeugt werden, sondern auf explizit
formulierten Regeln oder Wissensstrukturen beruhen. Die Adaptionslogik
wird dabei durch Expertinnen und Experten vorgegeben und vom System zur
Auswahl geeigneter Aktionen angewendet.

Typischerweise kommen If-Then-Regeln beziehungsweise
Condition-Action-Regeln zum Einsatz. Auf Grundlage beobachteter Merkmale
eines Lernenden werden dabei bestimmte Inhalte, Navigationsmöglichkeiten
oder Unterstützungsmaßnahmen aktiviert. Ein einfaches Beispiel wäre die
Regel, einem Lernenden bei wiederholten Fehlern zusätzliche Erklärungen
oder Übungsaufgaben bereitzustellen. Darüber hinaus können auch
symbolische Wissensrepräsentationen und regelbasierte
Inferenzmechanismen zur Ableitung adaptiver Entscheidungen genutzt
werden \citep{Wang2025}.

Ein wesentlicher Vorteil dieser Verfahren liegt in ihrer hohen
Transparenz und Nachvollziehbarkeit. Da die Entscheidungsregeln explizit
definiert sind, lässt sich die Ursache einer bestimmten
Adaptionsentscheidung unmittelbar nachvollziehen. Gleichzeitig sind
regelbasierte Verfahren häufig stark von Expertenwissen abhängig und
können neue oder unerwartete Situationen nur eingeschränkt
berücksichtigen. Im Vergleich zu datengetriebenen Verfahren besitzen sie
daher oftmals eine geringere Flexibilität und Anpassungsfähigkeit.

Innerhalb der vorgeschlagenen Klassifikation werden regel- und
wissensbasierte Verfahren als eigenständige Unterklasse geführt, da die
Adaptionslogik hier nicht gelernt oder optimiert wird, sondern explizit
modelliert ist. Typische Vertreter dieser Kategorie sind regelbasierte
Tutorensysteme, Expertensysteme sowie Condition-Action-basierte
Adaptionsmechanismen.

\subsubsection{Informationsmaximierende
Selektionsverfahren}\label{informationsmaximierende-selektionsverfahren}

Informationsmaximierende Selektionsverfahren bilden die zweite
Unterklasse der Adaptionsentscheidung. Im Mittelpunkt dieser Verfahren
steht die Auswahl derjenigen Aufgabe oder Testfrage, die auf Grundlage
des aktuell geschätzten Lernendenzustands den größten diagnostischen
Informationsgewinn verspricht. Im Unterschied zu Verfahren der
Zustandsmodellierung besteht ihre unmittelbare Ausgabe somit nicht in
einer Fähigkeitsschätzung, sondern in der Entscheidung, welches Item als
Nächstes präsentiert werden soll.

Besonders verbreitet sind diese Verfahren im Bereich des Computerized
Adaptive Testing (CAT). Dort wird die aktuelle Fähigkeitsschätzung eines
Lernenden fortlaufend genutzt, um die nächste Aufgabe adaptiv
auszuwählen. Ziel dieser Verfahren ist es, auf Basis der aktuellen
Fähigkeitsschätzung jene Aufgabe auszuwählen, die den größten
zusätzlichen Informationsgewinn für die weitere Schätzung des
Fähigkeitsniveaus liefert. Die Aufgabenauswahl erfolgt dabei
typischerweise anhand informationsbasierter Auswahlkriterien, die
bestimmen, welche Aufgabe unter den verfügbaren Kandidaten den größten
zusätzlichen Erkenntnisgewinn liefert \citep{Liu2024}.

Zu den wichtigsten Verfahren dieser Unterklasse zählen
informationsmaximierende Item-Selection-Kriterien wie die
Maximum-Fisher-Information-Auswahl oder Kullback-Leibler-basierte
Selektionsverfahren. Darüber hinaus spielen Verfahren des Content
Balancing sowie Mechanismen zur Kontrolle der Itemexposition eine
wichtige Rolle, um in realen Testsystemen sowohl diagnostische Qualität
als auch faire Nutzung der Aufgabenbanken sicherzustellen
\citep{Liu2024}.

Charakteristisch für diese Modellfamilie ist die explizite Optimierung
der Aufgabenauswahl unter diagnostischen und testtheoretischen
Randbedingungen. Die Verfahren wählen Aufgaben nicht primär aufgrund
inhaltlicher Ähnlichkeit oder Präferenzen aus, sondern anhand ihres
erwarteten Informationsbeitrags zur Fähigkeitsmessung. Dadurch
unterscheiden sie sich grundlegend von Recommender-Systemen oder
Lernpfadverfahren, deren Ziel stärker auf Personalisierung und
Inhaltsauswahl ausgerichtet ist.

Innerhalb der vorgeschlagenen Klassifikation werden
informationsmaximierende Selektionsverfahren als eigenständige
Unterklasse geführt, da die Adaptionsentscheidung hier explizit als
Optimierungsproblem der diagnostischen Informationsgewinnung formuliert
wird.

\subsubsection{Ressourcenempfehlung}\label{ressourcenempfehlung}

Verfahren der Ressourcenempfehlung bilden die dritte Unterklasse der
Adaptionsentscheidung. Im Mittelpunkt dieser Verfahren steht die Auswahl
geeigneter Lernressourcen, Aufgaben oder Inhalte auf Grundlage
verfügbarer Informationen über den Lernenden. Im Unterschied zu
informationsmaximierenden Selektionsverfahren erfolgt die Entscheidung
dabei nicht primär nach diagnostischem Informationsgewinn, sondern nach
der erwarteten Relevanz, Passung oder Nützlichkeit eines Lernobjekts für
einen bestimmten Lernenden.

Besonders verbreitet sind diese Verfahren im Bereich der Educational
Recommender Systems (ERS). Ziel solcher Systeme ist es, Lernenden
diejenigen Ressourcen vorzuschlagen, die ihren Interessen, ihrem
Wissensstand, ihren Lernzielen oder ihren bisherigen Interaktionen
möglichst gut entsprechen. Die Literatur unterscheidet dabei
insbesondere zwischen content-basierten, kollaborativen,
wissensbasierten sowie hybriden Empfehlungsverfahren
\citep{daSilva2023}.

Content-basierte Verfahren empfehlen Lernressourcen auf Grundlage ihrer
Eigenschaften und Merkmale. Dabei werden Inhalte ausgewählt, die
Ähnlichkeiten zu bereits genutzten oder positiv bewerteten Ressourcen
aufweisen. Kollaborative Verfahren nutzen dagegen Präferenz- oder
Interaktionsmuster mehrerer Lernender und gehen davon aus, dass Personen
mit ähnlichem Verhalten auch ähnliche Ressourcen bevorzugen. Wissens-
beziehungsweise ontologiebasierte Verfahren greifen auf explizit
modelliertes Domänenwissen, Lernziele oder semantische Beziehungen
zwischen Lernobjekten zurück, um geeignete Empfehlungen abzuleiten
\citep{daSilva2023}.

In der aktuellen Literatur besitzen hybride Empfehlungsverfahren eine
besondere Bedeutung. Sie kombinieren mehrere Empfehlungsstrategien,
beispielsweise content-basierte und kollaborative Verfahren, um die
jeweiligen Stärken verschiedener Ansätze miteinander zu verbinden und
deren Einschränkungen teilweise zu reduzieren. Insbesondere Probleme wie
der Cold-Start-Effekt, bei dem für neue Lernende oder Lernressourcen
zunächst kaum Informationen vorliegen, sowie Datensparsamkeit (Data
Sparsity), also eine geringe Anzahl verfügbarer Interaktionsdaten,
können durch die Kombination mehrerer Empfehlungsstrategien teilweise
reduziert werden. Typische Vertreter dieser Unterklasse sind
kollaborative Filterverfahren, content-basierte Empfehlungssysteme,
wissens- bzw. ontologiebasierte Recommender sowie hybride
Empfehlungssysteme \citep{daSilva2023}.

Charakteristisch für diese Modellfamilie ist die Verbindung von
Zustandsinformationen über den Lernenden mit der Auswahl konkreter
Lernressourcen. Die Verfahren bilden damit die Schnittstelle zwischen
Learner Modeling und inhaltlicher Personalisierung. Während
Zustandsmodelle Informationen über Wissen, Fähigkeiten oder Präferenzen
bereitstellen, entscheiden Empfehlungssysteme darüber, welche Inhalte
auf Grundlage dieser Informationen als Nächstes vorgeschlagen werden
sollen.

Innerhalb der vorgeschlagenen Klassifikation werden Verfahren der
Ressourcenempfehlung als eigenständige Unterklasse geführt, da ihre
primäre Aufgabe in der Auswahl geeigneter Lernobjekte auf Basis von
Relevanz-, Ähnlichkeits- oder Präferenzkriterien besteht. Typische
Vertreter dieser Kategorie sind content-basierte Recommender,
kollaborative Filterverfahren, wissensbasierte Empfehlungssysteme sowie
hybride Educational Recommender Systems.

\subsubsection{Lernpfadempfehlung und
Sequenzplanung}\label{lernpfadempfehlung-und-sequenzplanung}

Verfahren der Lernpfadempfehlung und Sequenzplanung bilden die vierte
Unterklasse der Adaptionsentscheidung. Im Unterschied zu klassischen
Empfehlungssystemen steht hier nicht die Auswahl eines einzelnen
Lernobjekts im Mittelpunkt, sondern die Planung ganzer Sequenzen von
Lernaktivitäten. Ziel dieser Verfahren ist die Konstruktion
personalisierter Lernpfade, die den aktuellen Wissensstand, individuelle
Lernziele sowie weitere Randbedingungen eines Lernenden berücksichtigen.

Die Literatur beschreibt personalisierte Lernpfade als einen zentralen
Bestandteil adaptiver Lernsysteme. Während Recommender-Systeme
typischerweise einzelne Ressourcen empfehlen, versuchen
Lernpfadverfahren, mehrere Lernobjekte in eine sinnvolle Reihenfolge zu
bringen und dadurch den gesamten Lernprozess zu strukturieren. Nabizadeh
et al.~unterscheiden hierzu die Ansätze Course Generation, Sequential
Pattern Recognition und Course Sequence. Gemeinsam ist diesen Verfahren,
dass sie zeitliche und inhaltliche Abhängigkeiten zwischen
Lernaktivitäten berücksichtigen und Lernen als Sequenzierungs-
beziehungsweise Optimierungsproblem modellieren \citep{Nabizadeh2020}.

Auf algorithmischer Ebene kommen unterschiedliche Such-, Optimierungs-
und Planungsverfahren zum Einsatz. Nabizadeh et al.~nennen insbesondere
graphbasierte Verfahren, Constraint-based Approaches,
Case-Based-Reasoning-Ansätze sowie evolutionäre und schwarmbasierte
Optimierungsverfahren wie Genetic Algorithms, Ant Colony Optimization
und Particle Swarm Optimization \citep{Nabizadeh2020}. Gemeinsam ist
diesen Verfahren, dass sie nicht lediglich einzelne Lernobjekte
auswählen, sondern ganze Lernsequenzen als Optimierungsproblem
modellieren.

Charakteristisch für diese Modellfamilie ist die explizite
Berücksichtigung zeitlicher und inhaltlicher Abhängigkeiten zwischen
Lernobjekten. Die Verfahren modellieren Lernen somit nicht als Folge
isolierter Empfehlungen, sondern als zusammenhängenden Prozess, dessen
einzelne Schritte aufeinander aufbauen.

Innerhalb der vorgeschlagenen Klassifikation werden Verfahren der
Lernpfadempfehlung und Sequenzplanung als eigenständige Unterklasse
geführt, da ihre primäre Aufgabe in der Planung vollständiger
Lernsequenzen besteht. Typische Vertreter dieser Kategorie sind
Course-Generation-Verfahren, Sequential-Pattern-Mining-Ansätze,
graphbasierte Lernpfadverfahren sowie Optimierungsverfahren zur
Curriculum- und Lernpfadplanung.

\subsubsection{Sequentielle
Entscheidungsoptimierung}\label{sequentielle-entscheidungsoptimierung}

Verfahren der sequentiellen Entscheidungsoptimierung bilden die fünfte
Unterklasse der Adaptionsentscheidung. Im Mittelpunkt dieser Verfahren
steht nicht die Auswahl einzelner Lernressourcen oder die Planung
vorgegebener Lernpfade, sondern die Optimierung von
Entscheidungsstrategien über eine Folge von Interaktionen hinweg. Ziel
ist es, adaptives Verhalten zu erlernen, das langfristig den Lernerfolg
maximiert.

Charakteristisch für diese Modellfamilie ist die Betrachtung adaptiver
Entscheidungen als sequentielles Entscheidungsproblem. Die Auswahl einer
Lernaktivität beeinflusst dabei nicht nur den unmittelbaren Lernerfolg,
sondern auch zukünftige Zustände und Handlungsmöglichkeiten. Verfahren
dieser Klasse müssen daher zwischen Exploration und Exploitation
abwägen. Exploration bezeichnet das Ausprobieren bislang wenig genutzter
Aktionen, um zusätzliche Informationen zu gewinnen, während Exploitation
die Nutzung bereits bekannter erfolgreicher Strategien beschreibt
\citep{Riedmann2025}.

Einen frühen Ansatz zur Modellierung adaptiver Bildungsentscheidungen
stellen kontextuelle Bandits dar. Hier wird auf Grundlage des aktuellen
Lernendenzustands diejenige Aktion ausgewählt, die den erwarteten
kurzfristigen Lernerfolg maximiert. Im Gegensatz zu vollständigen
Reinforcement-Learning-Verfahren berücksichtigen Bandit-Modelle dabei
typischerweise keine langfristigen Auswirkungen späterer Entscheidungen
\citep{Riedmann2025}.

Den wichtigsten Forschungsbereich dieser Unterklasse bildet das
Reinforcement Learning. Dabei lernt ein Agent durch Interaktion mit der
Lernumgebung eine Strategie (Policy), die langfristige Belohnungen
maximiert. Die Literatur beschreibt insbesondere wertbasierte Verfahren
wie Q-Learning und Deep Q-Networks (DQN) sowie Policy-Gradient- und
weitere Policy-Learning-Verfahren als häufig eingesetzte Ansätze im
Bildungsbereich \citep{Riedmann2025}.

Im Unterschied zu klassischen Recommender-Systemen erfolgt die Auswahl
von Lernaktivitäten hier nicht primär auf Grundlage von Ähnlichkeits-,
Präferenz- oder Relevanzmaßen. Stattdessen wird die gesamte Folge von
Zuständen, Aktionen und Belohnungen berücksichtigt. Die
Adaptionsentscheidung wird somit als Optimierungsproblem über längere
Interaktionssequenzen modelliert.

Innerhalb der vorgeschlagenen Klassifikation werden Verfahren der
sequentiellen Entscheidungsoptimierung als eigenständige Unterklasse
geführt, da sie adaptive Entscheidungen explizit als Folgeproblem mit
Exploration, Exploitation und langfristiger Belohnungsmaximierung
modellieren. Typische Vertreter dieser Kategorie sind kontextuelle
Bandits, Q-Learning, Deep Q-Networks (DQN), Proximal Policy Optimization
(PPO) \citep{Riedmann2025}.

\section{Ergänzende
Klassifikationsdimensionen}\label{erguxe4nzende-klassifikationsdimensionen}

Die bislang entwickelte Taxonomie klassifiziert adaptive Lernverfahren
primär nach ihrer funktionalen Rolle innerhalb adaptiver Lernsysteme.
Dabei wird zwischen Verfahren der Zustandsmodellierung und Verfahren der
Adaptionsentscheidung unterschieden. Diese Einteilung erlaubt eine
systematische Zuordnung der in der Literatur beschriebenen Verfahren,
erfasst jedoch nicht alle relevanten Unterschiede zwischen einzelnen
Ansätzen.

Aus analytischer Sicht ist es daher sinnvoll, die funktionale
Klassifikation durch zusätzliche Querschnittsdimensionen zu ergänzen.
Diese beschreiben Eigenschaften, die sich über mehrere Kategorien hinweg
erstrecken und einen Vergleich unterschiedlicher Verfahren ermöglichen.
Insbesondere die verwendete Datenbasis, die zugrunde liegende
Modellierungsfamilie sowie die Interpretierbarkeit eines Verfahrens
werden in der Literatur wiederholt als zentrale Unterscheidungsmerkmale
hervorgehoben.

\subsection{Datenbasis}\label{datenbasis}

Die Datenbasis beschreibt, auf welchen Informationen ein Verfahren seine
Modellierung oder Entscheidungsfindung aufbaut. Adaptive Lernsysteme
nutzen hierzu eine Vielzahl unterschiedlicher Datentypen. Dazu zählen
insbesondere Leistungsdaten wie richtige oder falsche Antworten,
Interaktions- und Prozessdaten wie Bearbeitungszeiten, Navigationspfade
oder Hint-Nutzung, Strukturwissen über Domänen und Lernobjekte sowie
Profil-, Präferenz- und Kontextinformationen der Lernenden.

Die Datenbasis stellt eine wichtige ergänzende Klassifikationsdimension
dar, da Verfahren derselben funktionalen Kategorie auf unterschiedlichen
Informationsquellen beruhen können. So zählen Rasch-Modelle, Bayesian
Knowledge Tracing und Deep Knowledge Tracing trotz ihrer gemeinsamen
Zuordnung zur Zustandsmodellierung zu unterschiedlichen Verfahrenstypen
und nutzen verfügbare Lerndaten in unterschiedlicher Form. Während
psychometrische Verfahren überwiegend Antwortdaten verarbeiten, beziehen
tief-neuronale Ansätze häufig längere Interaktionssequenzen ein. Analog
unterscheiden sich Verfahren der Adaptionsentscheidung hinsichtlich der
verwendeten Informationsquellen, etwa zwischen Präferenzdaten in
Recommender-Systemen und Zustands-Aktions-Belohnungsdaten im
Reinforcement Learning.

\subsection{Modellierungsfamilie}\label{modellierungsfamilie}

Die Modellierungsfamilie beschreibt das zugrunde liegende mathematische
oder informatische Repräsentationsprinzip eines Verfahrens. Verfahren
mit ähnlicher Funktion können dabei auf sehr unterschiedlichen
Modellierungsansätzen beruhen. In der Literatur finden sich insbesondere
psychometrische, probabilistische, logistische, symbolische,
graphbasierte, faktorisierungsbasierte, metaheuristische sowie
tief-neuronale Modellierungsansätze.

Die Modellierungsfamilie verdeutlicht, dass Verfahren mit gleicher
funktionaler Rolle auf unterschiedlichen mathematischen Grundlagen
beruhen können. Innerhalb der Zustandsmodellierung finden sich
beispielsweise psychometrische Verfahren wie IRT-Modelle,
probabilistische Ansätze wie BKT, logistische Verfahren wie PFA sowie
tief-neuronale Verfahren wie DKT. Die funktionale Einordnung allein
erfasst diese Unterschiede nicht, weshalb die Modellierungsfamilie eine
sinnvolle ergänzende Beschreibungsebene darstellt.

\subsection{Interpretierbarkeit}\label{interpretierbarkeit}

Die Interpretierbarkeit beschreibt, in welchem Umfang die Entscheidungen
oder Vorhersagen eines Verfahrens für Menschen nachvollziehbar sind.
Diese Eigenschaft besitzt im Bildungsbereich eine besondere Bedeutung,
da adaptive Entscheidungen häufig gegenüber Lernenden, Lehrenden oder
Institutionen begründet werden müssen.

Besonders deutlich wird dies innerhalb der Zustandsmodellierung.
Psychometrische Modelle, Bayes-Netz-Ansätze oder regelbasierte Verfahren
gelten häufig als vergleichsweise gut nachvollziehbar, da ihre
Modellannahmen und Entscheidungsregeln explizit beschrieben werden
können. Demgegenüber erreichen tiefe neuronale Verfahren oftmals höhere
Vorhersageleistungen, erzeugen ihre Ergebnisse jedoch auf Grundlage
schwer nachvollziehbarer interner Repräsentationen. Die
Interpretierbarkeit stellt daher keine eigenständige
Klassifikationsachse dar, ergänzt die funktionale Taxonomie jedoch um
eine für den Bildungsbereich wichtige Bewertungsperspektive.

Die Interpretierbarkeit bildet daher eine wichtige ergänzende
Klassifikationsdimension, da sie einen zentralen Zielkonflikt adaptiver
Lernsysteme sichtbar macht: den Ausgleich zwischen Vorhersage-
beziehungsweise Optimierungsleistung einerseits und Transparenz sowie
Erklärbarkeit andererseits.

Die beschriebenen Dimensionen stellen keine eigenständigen
Klassifikationsachsen der vorgeschlagenen Taxonomie dar, ermöglichen
jedoch eine differenziertere Beschreibung und den Vergleich von
Verfahren innerhalb derselben funktionalen Kategorie.

\section{Diskussion und Grenzen der
Klassifikation}\label{diskussion-und-grenzen-der-klassifikation}

Die vorgeschlagene Taxonomie ist als synthetische Ordnung adaptiver
Lernverfahren zu verstehen. Sie reproduziert keine einzelne, in der
Literatur etablierte Gesamtklassifikation, sondern führt mehrere in
unterschiedlichen Forschungsbereichen entwickelte Teilordnungen zu einem
gemeinsamen Schema zusammen. Ihre Stärke liegt damit in der
Zusammenführung fragmentierter Literaturstränge aus Bereichen wie
Learner Modeling, Knowledge Tracing, Computerized Adaptive Testing,
Educational Recommender Systems und Reinforcement Learning. Gleichzeitig
stellt genau diese Synthese auch eine Grenze der Taxonomie dar.

Eine erste Einschränkung ergibt sich aus der Tatsache, dass die
vorgeschlagenen Kategorien nicht vollständig trennscharf sind. Besonders
deutlich zeigt sich dies bei hybriden Systemen, die in der Literatur
eher die Regel als die Ausnahme darstellen. Adaptive Lernumgebungen
kombinieren häufig mehrere Modellierungs- und Entscheidungsverfahren
miteinander. So können Knowledge-Tracing-Modelle als Grundlage für
Empfehlungssysteme dienen, Lernpfadverfahren auf Zustandsmodellen
aufbauen oder Reinforcement-Learning-Verfahren Informationen aus
psychometrischen oder probabilistischen Lernendenmodellen verwenden.

Einzelne Verfahren können daher je nach Betrachtungsperspektive mehreren
Kategorien zugeordnet werden. So lassen sich Bayes-Netz-basierte Ansätze
sowohl zur Modellierung von Lernzuständen als auch als Bestandteil von
Adaptionsentscheidungen einsetzen. Ebenso können Recommender-Systeme
Komponenten des Reinforcement Learning integrieren oder Elo-basierte
Verfahren Elemente psychometrischer Fähigkeitsmodelle aufweisen. Die
Kategorien sind daher als analytische Hauptzuordnungen und nicht als
gegenseitig ausschließende Klassen zu verstehen. Eine ausschließlich auf
gegenseitig ausschließenden Kategorien beruhende Taxonomie würde die in
der Praxis verbreitete Hybridität adaptiver Lernsysteme nur unzureichend
abbilden.

Für die vorgeschlagene Klassifikation werden hybride Systeme daher nach
der primären funktionalen Rolle des jeweiligen Moduls eingeordnet.
Zusätzliche Verfahren oder Modellierungsansätze werden über die
ergänzenden Klassifikationsdimensionen berücksichtigt. Ein System, das
beispielsweise Bayesian Knowledge Tracing zur Wissensschätzung und
anschließend ein Recommender-System zur Auswahl geeigneter
Lernressourcen verwendet, umfasst somit sowohl ein Verfahren der
Zustandsmodellierung als auch ein Verfahren der Adaptionsentscheidung.
Die Taxonomie versteht sich daher nicht als Zuordnung vollständiger
Systeme zu exakt einer Kategorie, sondern als Einordnung der jeweils
verwendeten Verfahren und Module.

Eine zweite Grenze betrifft die Dynamik des Forschungsfeldes. Adaptive
Lernsysteme entwickeln sich kontinuierlich weiter, insbesondere im
Bereich datengetriebener und neuronaler Verfahren. Neue
Modellarchitekturen oder hybride Ansätze können bestehende Kategorien
miteinander verbinden oder neue Untergruppen erforderlich machen. Die
vorgeschlagene Taxonomie ist daher nicht als abschließende Ordnung zu
verstehen, sondern als strukturierender Bezugsrahmen für den aktuellen
Forschungsstand.

Darüber hinaus weist die Literatur wiederholt auf Defizite in der
Evaluation adaptiver Lernverfahren hin. Mehrere Reviews stellen fest,
dass die Auswahl bestimmter Adaptionsverfahren häufig nur begrenzt
theoretisch oder empirisch begründet wird. Gleichzeitig konzentrieren
sich viele Studien auf technische Leistungsmaße wie
Vorhersagegenauigkeit oder Empfehlungsqualität, während pädagogische
Wirkungen, Lernerfolge oder langfristige Nutzungseffekte deutlich
seltener untersucht werden. Auch im Bereich des Reinforcement Learning
basiert ein erheblicher Teil der vorhandenen Evidenz auf Simulationen
statt auf realen Lernumgebungen. Die Aussagekraft von Vergleichen
zwischen Verfahren ist dadurch teilweise eingeschränkt.

Die ergänzenden Klassifikationsdimensionen Datenbasis,
Modellierungsfamilie und Interpretierbarkeit wurden eingeführt, um
einige dieser Einschränkungen abzumildern. Sie verdeutlichen, dass
Verfahren mit gleicher funktionaler Rolle auf unterschiedlichen
Datenquellen, Modellierungsansätzen oder Transparenzniveaus beruhen
können. Gleichzeitig ersetzen diese Dimensionen nicht die funktionale
Hauptstruktur der Taxonomie, sondern dienen als ergänzende
Beschreibungsperspektiven.

Insgesamt erscheint die Unterscheidung zwischen Zustandsmodellierung und
Adaptionsentscheidung dennoch als ein sachlich klarer und informatisch
anschlussfähiger Ausgangspunkt für die Klassifikation adaptiver
Lernverfahren. Innerhalb dieser beiden Hauptklassen lassen sich die in
der Literatur beschriebenen Verfahren konsistent einordnen, ohne die
Vielfalt der zugrunde liegenden Modellierungsansätze vollständig zu
verlieren.
